\chapter{Especificación de requisitos}
\label{ch:especificación de requisitos}

\paragraph*{}En este capítulo se llevará a cabo un análisis más detallado y técnico de los requisitos que componen el sistema.
Partiendo de los requisitos ya definidos en el TFG anterior, este trabajo mantiene las funcionalidades generales básicas del sistema, como la carga, limpieza y visualización de datos, así como la interacción con el asistente de inteligencia artificial. No obstante, los requisitos específicos asociados a los modelos clásicos ya implementados en el sistema base no se redefinen de forma individual, puesto que forman parte de la solución previa. A partir de dicha base bien fundamentada, este TFG incorpora nuevos requisitos orientados a la ampliación funcional y a la mejora de la capacidad analítica del dashboard.
\section{Requisitos funcionales}
\begin{itemize}
    \item \textbf{RF1.} El sistema debe permitir la interacción del usuario con el asistente de inteligencia artificial integrado.
    
    \item \textbf{RF2.} El sistema debe incorporar soporte multilingüe en la interfaz, permitiendo al menos el uso en castellano e inglés.
    
    \item \textbf{RF3.} El sistema debe permitir trabajar con nuevas variables del conjunto de datos, además de las ya contempladas en el sistema anterior.
    
    \item \textbf{RF4.} El sistema debe admitir el análisis de variables no binarias dentro del flujo de estudio.
    
    \item \textbf{RF5.} El sistema debe ampliar las visualizaciones disponibles para mejorar la interpretación de los resultados en los métodos de regresión de Cox y Test de Log-Rank.
    
    \item \textbf{RF6.} El sistema debe incorporar nuevas técnicas de análisis de supervivencia que amplíen la capacidad analítica del dashboard, cada una con sus respectivas interpretaciones.
    
    \item \textbf{RF7.} El sistema debe incorporar el modelo exponencial como nueva técnica de análisis de supervivencia.
    
    \item \textbf{RF8.} El sistema debe incorporar el modelo de Weibull como nueva técnica de análisis de supervivencia.
    
    \item \textbf{RF9.} El sistema debe incorporar Random Survival Forest como nueva técnica de análisis de supervivencia.
    
    \item \textbf{RF10.} El sistema debe permitir la comparación de resultados entre las distintas técnicas de análisis implementadas.
    
    \item \textbf{RF11.} El sistema debe incorporar una nueva solución de inteligencia artificial más rápida que la utilizada en el TFG anterior, manteniendo la generación de interpretaciones sobre los resultados estadísticos.
    
    \item \textbf{RF12.} El sistema debe mejorar la calidad de las explicaciones generadas por la inteligencia artificial a partir de los resultados estadísticos obtenidos.
    
    \item \textbf{RF13.} El sistema debe permitir la exportación de resultados, gráficas, tablas e informes en formato PDF para su uso en documentación o análisis posteriores.
\end{itemize}



\section{Requisitos no funcionales}
\begin{itemize}
    \item \textbf{RNF1.} El sistema debe ofrecer tiempos de respuesta adecuados durante la carga, limpieza y visualización de los datos.
    
    \item \textbf{RNF2.} El sistema debe ejecutar los análisis en un tiempo razonable, incluso tras la incorporación de nuevas técnicas de análisis de supervivencia.
    
    \item \textbf{RNF3.} El sistema debe reducir el tiempo de respuesta del módulo de inteligencia artificial con respecto a la solución utilizada en el TFG anterior, teniendo en cuenta métricas como el tiempo de respuesta, el tiempo de generación y el consumo de recursos.
    
    \item \textbf{RNF4.} El sistema debe garantizar la confidencialidad de los datos académicos utilizados en el análisis.
    
    \item \textbf{RNF5.} El sistema debe validar los archivos cargados para evitar la incorporación de contenido malicioso o no válido.
    
    \item \textbf{RNF6.} El sistema debe ofrecer una interfaz clara, comprensible y fácil de utilizar para facilitar la navegación y la comprensión de la información mostrada.
    
    \item \textbf{RNF7.} El sistema debe mostrar correctamente las visualizaciones y resultados en los navegadores de uso habitual.
    
    \item \textbf{RNF8.} El sistema debe mantener la estabilidad de las funcionalidades ya implementadas en la versión previa.
    
    \item \textbf{RNF9.} El sistema debe asegurar la coherencia entre los resultados estadísticos obtenidos, las gráficas generadas y las explicaciones producidas por la inteligencia artificial.
    
    \item \textbf{RNF10.} El sistema debe facilitar futuras ampliaciones funcionales sin necesidad de modificar por completo la estructura principal de la aplicación.
    
    \item \textbf{RNF11.} El sistema debe permitir un despliegue reproducible en entorno local con las dependencias necesarias correctamente definidas.
\end{itemize}

\section{Requisitos de información}
\begin{itemize}
    \item \textbf{RI1.} El sistema debe gestionar un conjunto de datos estructurado que contenga información académica de los estudiantes, incluyendo variables relevantes para el análisis de supervivencia.
    
    \item \textbf{RI2.} El sistema debe permitir la incorporación de nuevas variables al conjunto de datos, incluyendo variables no binarias.
    
    \item \textbf{RI3.} El sistema debe almacenar y procesar correctamente las variables necesarias para la aplicación de modelos de análisis de supervivencia.
    
    \item \textbf{RI3.1.} El sistema debe identificar correctamente las variables mínimas necesarias para el análisis, incluyendo un identificador del estudiante, una variable temporal y una variable que represente el evento o resultado observado.
    
    \item \textbf{RI3.2.} El sistema debe representar correctamente la información asociada a datos censurados, diferenciando los casos en los que el abandono se ha producido de aquellos en los que no se ha observado durante el periodo de estudio.
    
    \item \textbf{RI4.} El sistema debe generar y mantener los resultados derivados de los distintos modelos aplicados, incluyendo métricas, estimaciones y representaciones gráficas.
    
    \item \textbf{RI5.} El sistema debe proporcionar información interpretativa sobre los resultados obtenidos, facilitando su comprensión mediante el módulo de inteligencia artificial.
    
    \item \textbf{RI6.} El sistema debe permitir la exportación de la información generada, incluyendo resultados, tablas y gráficos, en formatos reutilizables.
\end{itemize}
